% =========================================================
% Lab 1 — Basic Statistics (solution)
% =========================================================
\documentclass[a4paper,11pt]{article}

\newcommand{\TPnumero}{Lab 1 -- Solution}
% =========================================================
% Tutorial / lab preamble — Time Series Analysis module
% article class + fancyhdr + tcolorbox + listings (English)
% Author : Aymen Ben Brik (aymen.benbrik@esprit.tn)
% =========================================================

\usepackage[utf8]{inputenc}
\usepackage[T1]{fontenc}
\usepackage[english]{babel}
\usepackage{lmodern}
\usepackage{amsmath, amssymb, amsthm, mathtools, bm}
\usepackage{graphicx}
\usepackage{booktabs}
\usepackage{array}
\usepackage{multirow}
\usepackage{colortbl}
\usepackage{xcolor}
\usepackage{tikz}
\usetikzlibrary{positioning, arrows.meta, calc, shapes, fit, backgrounds, matrix}
\usepackage{listings}
\usepackage[most]{tcolorbox}
\usepackage{geometry}
\geometry{a4paper, margin=2cm, headheight=15pt}
\usepackage{fancyhdr}
\usepackage[hidelinks]{hyperref}
\usepackage{enumitem}

% --- Colours ---
\definecolor{espritBlue}{RGB}{30, 60, 110}
\definecolor{espritAccent}{RGB}{200, 80, 30}
\definecolor{boxEx}{RGB}{30, 60, 110}
\definecolor{boxQ}{RGB}{180, 120, 0}
\definecolor{boxC}{RGB}{30, 100, 60}
\definecolor{boxAtt}{RGB}{200, 80, 30}

% --- Listings ---
\definecolor{codebg}{RGB}{248,248,248}
\definecolor{codeframe}{RGB}{220,220,220}
\definecolor{keyword}{RGB}{0,82,155}
\definecolor{comment}{RGB}{88,110,117}
\definecolor{string}{RGB}{133,153,0}

\lstdefinestyle{pythonstyle}{
  language=Python,
  basicstyle=\ttfamily\small,
  keywordstyle=\color{keyword}\bfseries,
  commentstyle=\color{comment}\itshape,
  stringstyle=\color{string},
  backgroundcolor=\color{codebg},
  frame=single,
  rulecolor=\color{codeframe},
  frameround=tttt,
  breaklines=true,
  showstringspaces=false,
  tabsize=2
}
\lstdefinestyle{Rstyle}{
  language=R,
  basicstyle=\ttfamily\small,
  keywordstyle=\color{keyword}\bfseries,
  commentstyle=\color{comment}\itshape,
  stringstyle=\color{string},
  backgroundcolor=\color{codebg},
  frame=single,
  rulecolor=\color{codeframe},
  frameround=tttt,
  breaklines=true,
  showstringspaces=false,
  tabsize=2
}
\lstset{style=pythonstyle}

% --- Common macros (configurable path) ---
\providecommand{\macrosPath}{../../preamble/macros.tex}
\input{\macrosPath}

% --- Exercise counter ---
\newcounter{excounter}
\renewcommand{\theexcounter}{\arabic{excounter}}

\newtcolorbox[use counter=excounter]{exercise}[1][]{
  enhanced, breakable,
  colback=boxEx!5, colframe=boxEx, fonttitle=\bfseries,
  title={Exercise~\theexcounter\ifstrempty{#1}{}{ -- #1}},
  arc=2pt, boxrule=0.6pt
}
\newtcolorbox{question}[1][]{
  enhanced, breakable,
  colback=boxQ!5, colframe=boxQ, fonttitle=\bfseries,
  title={Question\ifstrempty{#1}{}{ -- #1}},
  arc=2pt, boxrule=0.5pt
}
\newtcolorbox{solution}[1][]{
  enhanced, breakable,
  colback=boxC!5, colframe=boxC, fonttitle=\bfseries,
  title={Solution\ifstrempty{#1}{}{ -- #1}},
  arc=2pt, boxrule=0.5pt
}
\newtcolorbox{warning}[1][]{
  enhanced, breakable,
  colback=boxAtt!5, colframe=boxAtt, fonttitle=\bfseries,
  title={Warning\ifstrempty{#1}{}{ -- #1}},
  arc=2pt, boxrule=0.5pt
}

% --- Headers / footers ---
\pagestyle{fancy}
\fancyhf{}
\fancyhead[L]{\textsc{\moduleNom} -- \auteurNom}
\fancyhead[R]{\TPnumero}
\fancyfoot[C]{\thepage}
\fancyfoot[L]{\auteurInstitution}
\fancyfoot[R]{\auteurEmail}
\renewcommand{\headrulewidth}{0.4pt}
\renewcommand{\footrulewidth}{0.2pt}

% --- Placeholder ---
\providecommand{\TPnumero}{Lab}


\title{Lab 1: Basic statistics in Python\\
\large \emph{Reference solution}}
\author{\auteurNom\\\auteurEmail\\\auteurInstitution}
\date{\anneeUniversitaire}

\begin{document}
\maketitle

\noindent\emph{Reference Python code is in \texttt{correction.py}; this
document provides the analytical answers and expected numerical values.}

% =========================================================
\setcounter{excounter}{0}
\begin{exercise}[Empirical moments]\end{exercise}

\begin{solution}
\textbf{(1)} For $n = 10\,000$ Gaussian draws (seed 42):
\[
  \widehat\mu \approx 2.005, \quad
  \widehat\sigma^2 \approx 3.987, \quad
  \widehat S \approx -0.013, \quad
  \widehat K \approx 2.996.
\]
All four are close to their theoretical values $(2, 4, 0, 3)$, with a
discrepancy of order $O(1/\sqrt{n})$ for the mean and variance.

\textbf{(2)} For $n = 100$, the variability of $\widehat\mu$ is
$\sigma / \sqrt{100} = 0.2$, so a typical deviation from $\mu = 2$ is
of order $0.2$. The plot of $\widehat\mu$ vs $n$ shows convergence
toward $\mu$ at rate $1/\sqrt{n}$ (Law of Large Numbers).
\end{solution}

% =========================================================
\begin{exercise}[Method of Moments vs MLE for the Gaussian]\end{exercise}

\begin{solution}
\textbf{(1)} For the Gaussian, $\widehat\mu_{\text{MoM}} = \widehat\mu_{\text{MLE}}
= \bar X_n$ and $\widehat\sigma^2_{\text{MoM}} = \widehat\sigma^2_{\text{MLE}}
= \tfrac{1}{n}\sum(X_i - \bar X)^2$. The two functions return identical
values on the same sample.

\textbf{(2)} On $B = 1000$ Monte-Carlo replicates with $n = 20$,
$\sigma^2 = 4$:
\[
  \widehat{\mathrm{Var}}_{\text{MLE}}\text{ mean } \approx 0.95\, \sigma^2
  = 3.80,
  \qquad
  S^2_n\text{ mean } \approx \sigma^2 = 4.00.
\]
The factor $\tfrac{n - 1}{n} = 0.95$ is exactly the bias predicted by
theory.
\end{solution}

% =========================================================
\begin{exercise}[Sampling distribution of $\bar X_n$]\end{exercise}

\begin{solution}
\textbf{(1)--(2)} The empirical variance of $\bar X_n$ over the $B$
replicates is close to $\sigma^2 / n = 4/n$:

\noindent
\begin{tabular}{rcc}
\toprule
$n$ & $\sigma^2/n$ & empirical variance \\
\midrule
$10$   & $0.40$   & $\approx 0.40$ \\
$30$   & $0.133$  & $\approx 0.13$ \\
$100$  & $0.040$  & $\approx 0.04$ \\
$1000$ & $0.004$  & $\approx 0.004$ \\
\bottomrule
\end{tabular}

\medskip
The four histograms agree with the theoretical $\mathcal{N}(\mu,
\sigma^2/n)$ density.

\textbf{(3)} For $\mathrm{Exp}(1)$ ($\mu = \sigma^2 = 1$), the histogram
of $\bar X_n$ is visibly skewed for $n = 10$ but well-approximated by
$\mathcal{N}(1, 1/n)$ from $n \geq 30$ onwards — illustrating the CLT.
\end{solution}

% =========================================================
\begin{exercise}[CI and $t$-test on Atlanta temperatures]\end{exercise}

\begin{solution}
The dataset \texttt{AvTempAtlanta.txt} contains 138 years (1879--2016)
of monthly temperatures, stacked into $n = 138 \times 12 = 1656$
observations.

\textbf{(1)} Empirical statistics: $\bar X_n = 61.80$, $S_n = 13.37$.
The 95\% $t$-CI is
\[
  \mathrm{CI}_{0.95}(\mu)
  = 61.80 \pm t_{1655, 0.975} \cdot \tfrac{13.37}{\sqrt{1656}}
  \approx 61.80 \pm 0.64
  = [61.16,\, 62.45].
\]

\textbf{(2)} For $H_0 : \mu = 60$ vs $H_1 : \mu \neq 60$:
\[
  T = \frac{61.80 - 60}{13.37 / \sqrt{1656}} \approx 5.486,
  \qquad
  t_{1655, 0.975} \approx 1.961.
\]
$|T| \gg 1.961$, so we \textbf{reject} $H_0$ at the 5\% level. The
$p$-value is $\approx 4 \times 10^{-8}$. The mean monthly temperature is
significantly different from $60^\circ$F.

\textbf{(3)} \texttt{scipy.stats.ttest\_1samp(x, popmean=60)} returns
the same statistic and $p$-value (up to floating-point precision).
\end{solution}

% =========================================================
\begin{exercise}[Bonus: bootstrap CI]\end{exercise}

\begin{solution}
With $B = 5000$ resamples on the $n = 1656$ Atlanta observations:
\[
  \mathrm{CI}_{0.95}^{\text{boot}}(\mu) \approx [61.15,\, 62.45],
\]
essentially identical to the $t$-CI. Both methods agree because $n$ is
large (CLT very effective). The bootstrap would be more useful in a
small-sample, non-Gaussian setting where the $t$-approximation is
questionable.
\end{solution}

\end{document}
